% problem-000329 1 金矿氰化浸出
\section*{1. 金矿氰化浸出}
\textit{下列为分问。}

\subsection*{1-1}
写出⽤氰化物溶⾦反应和⽤Zn粉置换⾦的化学反应⽅程式。

\subsection*{1-2}
已知 $1 E ^ { \ominus } ( \mathrm { O } _ { 2 } / \mathrm { H } _ { 2 } \mathrm { O } ) = 1 . 2 2 9 \mathrm { V }$ $E ^ { \ominus } ( \mathrm { A u ^ { + } / A u } ) = 1 . 6 9 \mathrm { V }$ $K _ { a } ^ { \ominus } ( \mathrm { H C N } ) = 4 . 9 3 \times 1 0 ^ { - }$ −10， $\beta _ { 2 } ( \mathrm { [ A u ( C N ) _ { 2 } ] ^ { - } } ) = 2 . 0 0 \times$ $1 0 ^ { 3 8 }$ ， $F _ { \mathrm { \ell } } = 9 6 4 8 5 \mathrm { J \ V ^ { - 1 } \ m o l ^ { - 1 } \Omega _ { \mathrm { \ell } } }$ 。设配制的NaCN⽔溶液的浓度为 $1 . 0 0 { \times } 1 0 ^ { - 3 } \mathrm { m o l } \mathrm { L } ^ { - 1 }$ 、⽣成的 $\mathrm { \ A u ( C N ) _ { 2 } ] }$ −配离⼦的浓度为 $1 . 0 0 { \times } 1 0 ^ { - 4 } \mathrm { m o l } \mathrm { L } ^ { - }$ 1、空⽓中 $\mathrm { O _ { 2 } }$ 的体积分数为0.210，计算298 K时在空⽓中溶⾦反应的⾃由能变。

\subsection*{1-3}
当电极反应中有H+或OH−时，其电极电势�将受pH的影响，�-pH图体现了这种影响。 $E \mathrm { - p H }$ 图上有3种类型的线：电极反应的�与pH有关，为有⼀定斜率的直线；电极反应的�与pH⽆关，是⼀条平⾏于横坐标的直线；⾮氧化还原反应，是⼀条平⾏于纵坐标的直线。

电对的 $E \mathrm { - p H }$ 线的上⽅，是该电对的氧化型的稳定区， $E \mathrm { - p H }$ 线的下⽅，是还原型的稳定区；位于⾼位置线的氧化型易与低位置线的还原型反应；各曲线的交点所处的�和pH，是各电极的氧化型和还原型共存的条件。

下图是氰化法溶⾦过程的 $\mathrm { A u \mathrm { - C N \mathrm { - } H _ { 2 } O } }$ 系统的�-pH图，试借助该图对溶⾦反应和溶⾦的⼯艺条件进⾏讨$i \not \in \Sigma _ { \circ }$

（图：）
